\documentclass[a4paper]{article}     
\usepackage{amsfonts}
\usepackage{amsmath}
\usepackage{amssymb}
\usepackage{graphicx}

\newcommand{\Bgtan}{\operatorname{Bgtan}}

\begin{document}

\noindent \large Auteur: Mark Gijsbrechts \\
\noindent \large Studierichting: Informatica\\
\noindent \large Oefeningenreeks 2B nr. 1.2\\

\medskip

\normalsize



Druk $\sin(\Bgtan x)$ uit in x:\\

$\alpha = \Bgtan x$ met $\alpha \in \left] -\dfrac{\pi}{2}, \dfrac{\pi}{2} \right[$ \\

$\Leftrightarrow x = \tan \alpha$  \\

$x^2 = \tan^2 \alpha$ \\

$x^2 = \dfrac{\sin^2 \alpha}{\cos^2 \alpha}$ \\

$x^2 + 1 = \dfrac{\sin^2 \alpha}{\cos^2 \alpha} + \dfrac{\cos^2 \alpha}{\cos^2 \alpha}$ \\

$x^2 + 1 = \dfrac{\sin^2 \alpha + \cos^2 \alpha}{\cos^2 \alpha}$ \\

$x^2 + 1 = \dfrac{1}{\cos^2 \alpha}$ \\

$x^2 + 1 = \dfrac{1}{1 - \sin^2 \alpha}$ \\

$\dfrac{1}{x^2 + 1} = 1 - \sin^2 \alpha$ \\

$ \sin^2 \alpha = 1 - \dfrac{1}{x^2 + 1}$ \\

$ \sin^2 \alpha = \dfrac{(x^2 + 1) - 1}{x^2 + 1}$ \\

$ \sin^2 \alpha = \dfrac{x^2}{x^2 + 1}$ \\

$ \sin \alpha = \dfrac{x}{\sqrt{x^2 + 1}}$ \\

$ \sin(\Bgtan x) = \dfrac{x}{\sqrt{x^2 + 1}}$ \\

\begin{thebibliography}{999}
%\bibitem{a} Referentie
\end{thebibliography}
\end{document} 
